% -*-coding: utf-8 -*-

\def\usewhat{pdflatex}    % 你喜欢哪种编译方式，pdflatex dvipdfmx xelatex yap 
% 目前texlive2009和miktex2.8（ctex2.8基于miktex2.8）中均提供了ctex中文支持宏包，在使用pdflatex、dvipdfmx和xelatex时，
% 用户无需再进行中文配置，安装texlive09或miktex2.8后就可编译本模板。
% 对于windows用户，若使用xelatex编译方式，请根据XP及以前系统，还是Vista及以后系统，在setup/definition.tex文件中选择合适的字体，然后再编译。
% 由于dvipspdf编译方式相对其他编译方式比较麻烦，并且没有什么优势，texlive09中ctex宏包没有默认提供
% dvips的中文支持方式，本模板不再提供dvipspdf编译方式。

%定义xelatex的中间临时变量，若\usewhat为xelatex时，后面执行xelatx的相关选项
\def\atempxetex{xelatex} %这一项无需改动
%input "reference\reference.bib" %for winedt users
\def\version{1.9.4.20100419}         % 该变量仅用于模板文件的版本号控制，新的论文规范从1.9开始；
	% 自从版本1.9.4.20100419后开始提供基于texlive09，miktex2.8的TeX系统的支持，不再支持早期的tl08和miktex2.7及以前版本。
	% 请使用tl08和miktex2.7及以前版本的用户使用本模板的1.9.2.10090424版本。

\def \xuewei {Doctor}   % 定义学位 博士
%\def \xuewei {Master}    % 硕士

\def\oneortwoside{twoside} %定义单双面打印，只对硕士学位论文有效；
%\def\oneortwoside{oneside} % 硕士单面打印

\def\xueke{Engineering} % 定义学科 工学
%\def\xueke{Science}      % 理学
%\def\xueke{Management}   % 管理学
%\def\xueke{Arts}         % 艺术学

% -*-coding: utf-8 -*-

% 硕博类型 的一些定义

%% 导言区使用中文
\makeatletter
\@tempcnta=128
\loop \catcode\@tempcnta=13 \ifnum\@tempcnta<255 \advance \@tempcnta \@ne
\repeat
\makeatother

\newif\ifxueweidoctor %判断论文类型
\newif\ifxueweimaster
\def\temp{Doctor}
\ifx\temp\xuewei
  \xueweidoctortrue  \xueweimasterfalse
\fi
\def\temp{Master}
\ifx\temp\xuewei
  \xueweidoctorfalse  \xueweimastertrue
\fi

\ifxueweidoctor
  \newcommand{\cxuewei}{博士}
  \newcommand{\exuewei}{Doctor}
  \newcommand{\exueweier}{Doctoral}
  \newcommand{\xueweishort}{博}
\fi

\ifxueweimaster
  \newcommand{\cxuewei}{硕士}
  \newcommand{\exuewei}{Master}
  \newcommand{\exueweier}{Master}
  \newcommand{\xueweishort}{硕}
\fi


\ifxueweidoctor
  \def\oneortwoside{twoside}
\fi

\ifx\oneortwoside\undefined
  \def\oneortwoside{twoside}
\fi

\newif\ifoneortwoside
\def\temp{twoside}
\ifx\temp\oneortwoside
  \oneortwosidetrue
\else
  \oneortwosidefalse
\fi
    % 硕博类型

%下面的book选项中可以使用 draft 选项，使插入的图形只显示外框，以加快预览速度。
\documentclass[12pt,a4paper,openany,\oneortwoside]{book}
%\documentclass[12pt,a4paper,openright,\oneortwoside]{book}

% -*-coding: utf-8 -*-

% 图形支持宏包 为了使用pdftex 需要作相应判断
\usepackage{etex}%增加计数器总数（原来是256，宏包多，可能不够用），编译需基于eTeX
\usepackage{ifpdf}
%定义一个新判断命令 %有一个宏包ifpdf 可以完成这件事，应该比这个严谨
%\newif\ifpdf
%\ifx\pdfoutput\undefined
%   \pdffalse
%\else
%   \pdfoutput=1
%   \pdftrue
%\fi

%%%%%%%%彩色引用和书签%%%%%%
\ifx\atempxetex\usewhat 
	\usepackage[dvipdfm]{graphicx}
\else
	\ifpdf
		\usepackage[pdftex]{graphicx}
	\else
		\usepackage[dvips]{graphicx}
	\fi
\fi

\usepackage[%paperwidth=18.4cm, paperheight= 26cm,  % 版面控制宏包，定义规定的版面尺寸
            body={15true cm,22true cm},           %论文版芯145mm×210mm（包括页眉及页码则为15mm×240mm）
            bindingoffset=0 pt,                      %页码在版芯下边线之下隔行居中放置；
            %headheight=1.0true cm
            ]{geometry}
\usepackage{layouts}                    % 打印当前页面格式的宏包
\usepackage[sf]{titlesec}               % 控制标题的宏包
\usepackage{titletoc}                   % 控制目录的宏包
\usepackage[perpage,symbol]{footmisc}   % 脚注控制
\usepackage{fancyhdr}                   % fancyhdr宏包 页眉和页脚的相关定义
\usepackage{fancyref}
\usepackage{array}          %增强表格的功能 %可能与xeCJK宏包冲突，需放在xeCJK之前。

\ifx\atempxetex\usewhat
%\usepackage[slantfont,boldfont,CJKaddspaces]{xeCJK} %
%\CJKlanguage{zh-cn}
\usepackage{ctex}
\else
\usepackage{CJKutf8} %CJKpunct}   % 中文支持宏包
%\usepackage[UTF8]{ctex}
\usepackage{times}          % 使用Times字体的宏包
\fi

\usepackage{type1cm}        % tex1cm宏包，控制字体的大小
\usepackage{indentfirst}    % 首行缩进宏包
\usepackage{color}          % 支持彩色

\usepackage{amsmath}        % AMSLaTeX宏包 用来排出更加漂亮的公式
\usepackage{relsize}            % 调整公式字体大小 \mathsmaller \mathlarger
\usepackage{amssymb}
\usepackage{textcomp}   	  % 千分号等特殊符号
\usepackage{mathrsfs}       % 不同于\mathcal or \mathfrak 之类的英文花体字体
\usepackage{bm}              % 处理数学公式中的黑斜体的宏包
\usepackage[amsmath,thmmarks,hyperref]{ntheorem}% 定理类环境宏包，其中 amsmath 选项用来兼容 AMS LaTeX 的宏包
\usepackage{epsfig}         % eps图像
\usepackage[below]{placeins}%允许上一个section的浮动图形出现在下一个section的开始部分,还提供\FloatBarrier命令,使所有未处理的浮动图形立即被处理
%\usepackage{psfrag}        %替换eps图形中的文字
%\usepackage{floatflt}       % 图文混排用宏包,texlive2008默认不带此宏包，可以自己安装，若是不需要这个功能可以不管它。
\usepackage{rotating}       % 图形和表格的控制
%\usepackage{endfloat}      %可将浮动对象放置到文件的最后
\usepackage{setspace}       % 定制表格和图形的多行标题行距
\usepackage{flafter}       % 使得所有浮动体不能被放置在其浮动环境之前，以免浮动体在引述它的文本之前出现.
\usepackage{multirow}       %使用Multirow宏包，使得表格可以合并多个row格
\usepackage{booktabs}       % 表格，横的粗线；\specialrule{1pt}{0pt}{0pt}
\usepackage{longtable}      %支持跨页的表格。
%\usepackage[centerlast]{caption2}       %浮动图形和表格标题样式，这里已被ccaption完全替代。
\usepackage[hang,center]{subfigure}%支持子图 %centerlast 设置最后一行是否居中
\usepackage[subfigure]{ccaption} %,caption2,

%\usepackage{cite}          % 支持引用的宏包
\usepackage[sort&compress,numbers]{natbib}% 支持引用缩写的宏包
\usepackage{hypernat}

\usepackage{enumitem}       %使用enumitem宏包,改变列表项的格式
\usepackage{calc}           %长度可以用+ - * / 进行计算


% 生成有书签的pdf及其开关, 该宏包应放在所有宏包的最后, 宏包之间有冲突
\def\atemp{dvipspdf}\ifx\atemp\usewhat
\usepackage[dvips,unicode,
           bookmarksnumbered=true,
           bookmarksopen=true,
           colorlinks=false,    % 最后打印的时候可以改成false，这样字体都是黑色
           pdfborder={0 0 0},   % 去掉链接的边框
           citecolor=blue,
           linkcolor=red,
           anchorcolor=green,
           urlcolor=blue,            
           breaklinks=true
           ]{hyperref}
\usepackage{breakurl} %消除dvips的时候网页链接断行失效的问题。
\fi

\def\atemp{dvipdfmx}\ifx\atemp\usewhat
\usepackage[dvipdfm,unicode,           %dvi-->pdf 生成书签
            bookmarksnumbered=true,
            bookmarksopen=true,
            colorlinks=false,
            pdfborder={0 0 0},
            citecolor=blue,
            linkcolor=red,
            anchorcolor=green,
            urlcolor=blue,
            breaklinks=true
            ]{hyperref}
\fi

\def\atemp{pdflatex}\ifx\atemp\usewhat
\usepackage{cmap}                       %pdflatex编译时，可以生成可复制、粘贴的中文PDF文档
\usepackage[pdftex,unicode,
            %CJKbookmarks=true,
            bookmarksnumbered=true,
            bookmarksopen=true,
            colorlinks=false,
            pdfborder={0 0 0},
            citecolor=blue,
            linkcolor=red,
            anchorcolor=green,
            urlcolor=blue,           
            breaklinks=true
            ]{hyperref}
\fi

\def\atemp{yap}\ifx\atemp\usewhat
\usepackage[dvipdfmx,unicode,  %考虑到yap可以正反向搜索，dvipdf使yap中的链接有效。
            %CJKbookmarks=true,
            bookmarksnumbered=true,
            bookmarksopen=true,
            colorlinks=false,
            pdfborder={0 0 0},
            citecolor=blue,
            linkcolor=red,
            anchorcolor=green,
            urlcolor=blue,       
            breaklinks=true
            ]{hyperref}
\fi

\ifx\atempxetex\usewhat %\def\atempxetex{xelatex} main.tex中已定义;
\usepackage[xetex,           
            bookmarksnumbered=true,
            bookmarksopen=true,
            colorlinks=false,
            pdfborder={0 0 0},
            citecolor=blue,
            linkcolor=red,
            anchorcolor=green,
            urlcolor=blue,
            breaklinks=true,           
            naturalnames  %与algorithm2e宏包协调
            ]{hyperref}  
\fi

\usepackage{arydshln}       %分块矩阵画虚线，挺好用

\usepackage[boxed,linesnumbered,algochapter]{algorithm2e}  % 算法的宏包，注意宏包兼容性，先后顺序为float、hyperref、algorithm(2e)，否则无法生成算法列表
\usepackage{listings}
\lstset{numbers=left,frame=tb,basicstyle=\small,xleftmargin=47pt,xrightmargin=30pt,framexleftmargin=7mm,mathescape=true,emphstyle=\textbf}
\usepackage{threeparttable}
\usepackage[mediumspace,mediumqspace,Grey,squaren]{SIunits}
\definecolor{light-gray}{gray}{0.85}
 % 引用的宏包

% 论文包含的内容
% FIXME: I don't know why I have to put Resume last :( Anyway, the
% order here does not seem matter.
%\includeonly{
                %body/introduction,
                %body/conclusion
                %appendix/publications,
                %appendix/acknowledgements,
                %appendix/Authorization,
                %appendix/Resume
            %}
\graphicspath{{figures/}} %定义所有的eps文件在 figures 子目录下

\begin{document}
\ifx\atempxetex\usewhat\else
\begin{CJK*}{UTF8}{song}
\fi

% -*-coding: utf-8 -*-

%%%%%%%%%%%%%%%%%%%%%%%%%%%%%%%%%%%%%%%%%%%%%%%%%%%%%%%%%%%
% 重定义字体命令
% 注意win2000,没有 simsun,　最好到网上找一个
% 一些字体是office2000 带的
%%%%%%%%%%%%%%%%%%%%%%%%%%%%%%%%%%%%%%%%%%%%%%%%%%%%%%%%%%%
\ifx\atempxetex\usewhat 
  % xelatex调用系统字体；
  % 【请在命令行下用 "fc-list" 列出系统的字体，选择其对应的名字】
  % 【只列出中文字体的命令 "fc-list :lang=zh-cn" 】
\defaultfontfeatures{Mapping=tex-text}
\setmainfont{Times New Roman}   % 正文中的英文 采用 Times New Roman 字体；
\setsansfont{Times New Roman}   % 正文中的英文 采用 Times New Roman 字体；
\setCJKmainfont{SimSun}
\setCJKsansfont{SimHei}
\setCJKmonofont{LiSu}
\setCJKfamilyfont{song}{SimSun}
\setCJKfamilyfont{hei}{SimHei}
\setCJKfamilyfont{fs}{FangSong} % XP对应 FangSong_GB2312，Vista对应FangSong，注意根据系统切换
\setCJKfamilyfont{kai}{KaiTi}   % XP对应 KaiTi_GB2312，Vista对应KaiTi，注意根据系统切换
\setCJKfamilyfont{li}{LiSu}
\fi
\newcommand{\song}{\CJKfamily{song}}    % 宋体   (Windows自带simsun.ttf)
\newcommand{\fs}{\CJKfamily{fs}}        % 仿宋体 (Windows自带simfs.ttf)
\newcommand{\kai}{\CJKfamily{kai}}      % 楷体   (Windows自带simkai.ttf)
\newcommand{\hei}{\CJKfamily{hei}}      % 黑体   (Windows自带simhei.ttf)
\newcommand{\li}{\CJKfamily{li}}        % 隶书   (Windows自带simli.ttf)

%%%%%%%%%%%%%%%%%%%%%%%%%%%%%%%%%%%%%%%%%%%%%%%%%%%%%%%%%%%
% 重定义字号命令
%%%%%%%%%%%%%%%%%%%%%%%%%%%%%%%%%%%%%%%%%%%%%%%%%%%%%%%%%%%
\newcommand{\yihao}{\fontsize{26pt}{26pt}\selectfont}       % 一号, 1.倍行距
\newcommand{\xiaoyi}{\fontsize{24pt}{24pt}\selectfont}      % 小一, 1.倍行距
\newcommand{\erhao}{\fontsize{22pt}{22pt}\selectfont}       % 二号, 1.倍行距
\newcommand{\xiaoer}{\fontsize{18pt}{18pt}\selectfont}      % 小二, 单倍行距
\newcommand{\sanhao}{\fontsize{16pt}{16pt}\selectfont}      % 三号, 1.倍行距
\newcommand{\xiaosan}{\fontsize{15pt}{15pt}\selectfont}     % 小三, 1.倍行距
\newcommand{\sihao}{\fontsize{14pt}{14pt}\selectfont}       % 四号, 1.0倍行距
\newcommand{\banxiaosi}{\fontsize{13pt}{13pt}\selectfont}   % 半小四, 1.0倍行距
\newcommand{\xiaosi}{\fontsize{12pt}{12pt}\selectfont}      % 小四, 1.倍行距
\newcommand{\dawuhao}{\fontsize{11.5pt}{11.5pt}\selectfont} % 大五号, 单倍行距
\newcommand{\wuhao}{\fontsize{10.5pt}{10.5pt}\selectfont}   % 五号, 单倍行距
\newcommand{\xiaowu}{\fontsize{9.5pt}{9.5pt}\selectfont}    % 五号, 单倍行距
\newcommand{\banbanxiaosi}{\fontsize{12pt}{12pt}\selectfont}% 半半小四, 1.0倍行距

%避免宏包 hyperref 和 arydshln 不兼容带来的目录链接失效的问题。
\def\temp{\relax}
\let\temp\addcontentsline
\gdef\addcontentsline{\phantomsection\temp}
\newcommand*{\subfigencaptionlist}{} % 子图形加入目录时用

\makeatletter
\gdef\hitempty{}

\newcommand{\mr}[1]{\mathrm{#1}} %定义新命令，用\mr来代替\mathrm
\def \ReferenceEName {References} %%定义参考文献的标题
\def \ReferenceCName {参考文献}

%定义图表章节双标题命令
\newcommand{\figenname}{Fig.}
\newcommand{\listfigenname}{List of Figures}
\newfloatlist[chapter]{figen}{fen}{\listfigenname}{\figenname}
\newfixedcaption{\figencaption}{figen}
\renewcommand{\thefigen}{\thechapter-\arabic{figure}}
\renewcommand{\@cftmakefentitle}{\chapter*{\listfigenname\@mkboth{\bfseries\listfigenname}{\bfseries\listfigenname}}}


\newcommand{\FigureBiCaption}[2]
{\renewcommand{\figurename}{图}
\vspace{-1.6ex}%-0.5ex
\caption{\protect\setlength{\baselineskip}{1.5em}#1} %\protect\setlength{\baselinestretch}{1.3}\selectfont
%\vspace{-3.4mm}
%\figencaption{\protect\setlength{\baselineskip}{1.5em}#2}%
%\vspace{-3.4mm}
%%
%%其子图形加入目录
\def\hittemp{}
 \@for \hittemp:=\subfigencaptionlist \do {%
        \ifx \hitempty\hittemp\relax \else
          \addcontentsline{fen}{subfigen}{\protect\numberline\hittemp}
        \fi}
 \gdef\subfigencaptionlist{}
}

\setcounter{fendepth}{2} %英文图形目录的深度 1(只有一级目录) 2(有两级目录)
\setcounter{lofdepth}{2} %中文图形目录的深度 1(只有一级目录) 2(有两级目录)
\renewcommand*{\l@subfigure}{\@dottedxxxline{\ext@subfigure}{2}{3.8em}{1.5em}} %中文图形目录 subfigure
\gdef\l@subfigen{\@dottedtocline{2}{3.8em}{1.5em}}%英文图形目录 latex
\newif\ifsubfigtoc
\ifnum \tw@ > \@nameuse{c@fendepth} \subfigtocfalse \else \subfigtoctrue \fi
\newbox\tempbox
\renewcommand*{\subcapsize}{\wuhao} %设置子图英文标题的字号为五号；
\def\SubfigEnCaption{%
  \@ifnextchar [%
      {\SubfigEnCap}%
      {\SubfigEnCap[0pt]}
} \long\def\SubfigEnCap[#1]#2  %产生caption  有水平间距调整
{ \ifsubfigtoc %加入目录这个动作，一定要在 父图 之后,所在先暂存在 subfigencaptionlist
    \protected@xdef\subfigencaptionlist{\subfigencaptionlist,%
        {{\thesubfigure}\protect\ignorespaces{#2}}}
  \fi \vspace{1pt}
  \sbox{\tempbox}{\thesubfigure\hskip\subfiglabelskip #2}%
  \ifthenelse{\lengthtest{\wd\tempbox > \linewidth}}%
  {\\[-20pt]\hspace*{#1}\parbox[t]{\linewidth}{\flushleft\noindent\wuhao\selectfont\thesubfigure\hskip\subfiglabelskip \centering#2\hangafter=1\hangindent=15pt}}%
  {\\\hspace*{#1}\centerline{\wuhao\selectfont\thesubfigure\hskip\subfiglabelskip #2}}
} 

%\newcommand{\SubfigureCaption}[2]  % Two Parameters, the first one is the width of the subfigure,
%{
%\addtocounter{subfigure}{-1}       % the second one is the caption of the subfigure
%\vspace{-2ex}
%\subfigure[#2]{\rule{#1}{0pt}}
%}

\newcommand{\tblenname}{Table} %define tbl instead of table
\newcommand{\listtblenname}{List of Tables}
\newfloatlist[chapter]{tblen}{ten}{\listtblenname}{\tblenname}
\newfixedcaption{\tblencaption}{tblen}
\renewcommand{\thetblen}{\thechapter-\arabic{table}}% 将tblen换成table，因为table和tablen编号一致，而tablen在\longbitoccaption定义中无效。
\renewcommand{\@cftmaketentitle}{\chapter*{\listtblenname\@mkboth{\bfseries\listtblenname}{\bfseries\listtblenname}}}


\newcommand{\TableBiCaption}[2]
{
\renewcommand{\tablename}{表}
%\vspace{1.6ex}%-0.5ex
\caption{\protect\setlength{\baselineskip}{1.5em}#1}
%\vspace{-2ex}
%\tblencaption{\protect\setlength{\baselineskip}{1.5em}#2}
%\vspace{1ex}
}

%%%% 长表格的caption在中英文表格目录中正常显示
\setlength{\LTcapwidth}{\textwidth} %长表格宽度
\def\@cont@LT@LTBiToeCaption#1[#2]#3{%
  \LT@makecaption#1\fnum@table{#3}%
  \def\@tempa{#2}%
  \ifx\@tempa\@empty\else
    {\let\\\space
      %\phantomsection
      \addcontentsline{ten}{tblen}{\protect\numberline{\thetable}{#2}}}%%\addcontentsline{lot}{table}{\protect\numberline{}{#2}}}%
  \fi}
\def\LT@c@ption#1[#2]#3{%
  \LT@makecaption#1\fnum@table{#3}%
  \def\@tempa{#2}%
  \ifx\@tempa\@empty\else
     {\let\\\space
     %\phantomsection
     \addcontentsline{lot}{table}{\protect\numberline{\thetable}{#2}}}%
  \fi}
\let\@cont@oldLT@c@ption\LT@c@ption
\newcommand*{\LTBiTocCaption}[5]{
  \@if@contemptyarg{#1}{\caption{\centering #2}}{\caption[#1]{\centering #2}}%
  \global\let\@cont@oldtablename\tablename
  \gdef\tablename{Table} %#3
  \global\let\LT@c@ption\@cont@LT@LTBiToeCaption
  \\
  \@if@contemptyarg{#4}{\caption{\centering #5}}{\caption[#4]{\centering #5}}%
  \global\let\tablename\@cont@oldtablename
  \global\let\LT@c@ption\@cont@oldLT@c@ption}

\renewcommand{\cfttblendotsep}{1} %自定义图表目录中的点间距大小
\renewcommand{\cftfigendotsep}{1}

%\renewcommand{\tablename}{表}  %jdg提供的一种方法，英文长表会添加到中文目录中去，上面定义
%\newcommand{\LTBiCaption}[2]   %\bicaptiontwotoc 主要就是解决这个问题。
%{%
%\caption{#1} \gdef\tablename{Table}
%\\ %[-3.5ex]
%\caption{#2}
%\gdef\tablename{表}\\ %[-1.5ex]
%}

%%%---公式中符号描述----start----
%\begin{formulasymb}{式中}{-3pt}%-3pt,-20pt调与上方的间距。
%  \fdesfirst{第一标签}{控制控制控制控制控制}
%  \fdes{其他标签}{控制控制控制控制控制}
%\end{formulasymb}
\newenvironment{formulades}[1]%
{\noindent\begin{list}{}{%
\setlength\topsep{0pt}
\settowidth{\labelwidth}{#1}
\setlength{\labelsep}{1mm}
\setlength{\leftmargin}{\labelwidth+\labelsep}
}}{\end{list}}
\newenvironment{formulasymb}[2]%-\!-\!-\!-
{\vspace*{#2}\newcommand{\fdesfirst}[2]%
{\begin{formulades}{#1\hspace*{26pt}##1~\cdash}\item[#1\hspace*{26pt}##1~\cdash]{##2}\end{formulades}\vspace*{-21pt}}%自己调距
\newcommand{\fdes}[2]{\begin{formulades}{#1\hspace{26pt}##1~\cdash}\item[##1~\cdash]{##2}\end{formulades}\vspace*{-21pt}}}%自己调距
{\vspace{21pt}\relax}%21pt调距
%%----公式中符号描述----end-----

%% ---- 左对齐的公式 start-----   \begin{flualign} a=c \end{flualign}
\newenvironment{flualign}{%
    \@fleqntrue
    \@mathmargin = -1sp
    \@mathmargin\leftmargini minus\leftmargini
    \let\mathindent=\@mathmargin
  \start@align\@ne\st@rredfalse\m@ne
}{%
  \math@cr \black@\totwidth@
  \egroup
  \ifingather@
    \restorealignstate@
    \egroup
    \nonumber
    \ifnum0=`{\fi\iffalse}\fi
  \else
    $$%
  \fi
  \ignorespacesafterend
  \@fleqnfalse
}
%%  ---- 左对齐的公式  end----

%重新定义BiChapter命令，可实现标题手动换行，但不影响目录
\def\BiChapter{\relax\@ifnextchar [{\@BiChapter}{\@@BiChapter}}
\def\@BiChapter[#1]#2#3{\chapter[#1]{#2}
    \addcontentsline{toe}{chapter}{\bfseries \xiaosi Chapter \thechapter\hspace{0.5em} #3}}
\def\@@BiChapter#1#2{\chapter{#1}
    \addcontentsline{toe}{chapter}{\bfseries \xiaosi Chapter \thechapter\hspace{0.5em}{\boldmath #2}}}
%\newcommand{\BiChapter}[2]
%{
%    \chapter{#1}
%    \addcontentsline{toe}{chapter}{\bfseries Chapter \thechapter\hspace{0.5em} #2}
%}

\newcommand{\BiSection}[2]
{   \section{#1}
    \addcontentsline{toe}{section}{\protect\numberline{\csname thesection\endcsname}#2}
}

\newcommand{\BiSubsection}[2]
{    \subsection{#1}
    \addcontentsline{toe}{subsection}{\protect\numberline{\csname thesubsection\endcsname}#2}
}

\newcommand{\BiSubsubsection}[2]
{    \subsubsection{#1}
    \addcontentsline{toe}{subsubsection}{\protect\numberline{\csname thesubsubsection\endcsname}#2}
}

\newcommand{\BiAppendixChapter}[2] % 该附录命令适用于发表文章，简历等
{\phantomsection
\markboth{#1}{#1}%\markboth{\MakeUppercase{#1}}{\MakeUppercase{#1}}
\addcontentsline{toc}{chapter}{\hei #1}
\addcontentsline{toe}{chapter}{\bfseries \xiaosi #2}  \chapter*{#1}
}

\newcommand{\BiAppChapter}[2]    % 该附录命令适用于有章节的完整附录
{\phantomsection  \chapter{#1}   %\markboth{\MakeUppercase{#1}}{\MakeUppercase{#1}} %为了winedt中project tree中toc正确显示，不要挪到下一行；
%\addcontentsline{toc}{chapter}{\hei #1}
 \addcontentsline{toe}{chapter}{\bfseries \xiaosi Appendix \thechapter~~#2}
}

\renewcommand{\thefigure}{\arabic{chapter}-\arabic{figure}}%使图编号为 7-1 的格式 %\protect{~}
%\renewcommand\fnum@figure{\figurename\nobreakspace\thefigure\protect{~~~~~~~~~}} %

\renewcommand{\thesubfigure}{\alph{subfigure})}%使子图编号为 a)的格式
\renewcommand{\p@subfigure}{\thefigure ~~} %%使子图引用为 7-1 a) 的格式，母图编号和子图编号之间用~~加一个空格 
%\renewcommand{\thesubfigure}{\alph{subfigure}}
%\renewcommand{\p@subfigure}{\thefigure} %%使子图引用为 7-1a 的格式
%\renewcommand{\@thesubfigure}{\thesubfigure)\hskip\subfiglabelskip}%使子图编号为 a)的格式

\renewcommand{\thetable}{\arabic{chapter}-\arabic{table}}%%使表编号为 7-1 的格式
\renewcommand{\theequation}{\arabic{chapter}-\arabic{equation}}%%使公式编号为 7-1 的格式

%----------- 设置算法标题形式，由“Algorithm 2.1：算法标题” 改为“算法 2-1 算法标题”，并且定义算法的双标题命令；---------%

\newcommand{\algoenname}{Algo.} %算法英文标题
\newcommand{\listalgoenname}{List of Algorithms}
\newfloatlist[chapter]{algoen}{aen}{\listalgoenname}{\algoenname}
\newfixedcaption{\algoencaption}{algoen}
\renewcommand{\thealgoen}{\thechapter-\arabic{algocf}}
\renewcommand{\@cftmakeaentitle}{\chapter*{\listalgoenname\@mkboth{\bfseries\listalgoenname}{\bfseries\listalgoenname}}}

\renewcommand{\algorithmcfname}{算法}
\setlength\AlCapSkip{1.2ex}
\SetAlgoSkip{1pt}
\renewcommand{\algocf@captiontext}[2]{\wuhao#1\algocf@typo ~ \AlCapFnt{}#2} % text of caption
\expandafter\ifx\csname algocf@within\endcsname\relax% if \algocf@within doesn't exist
\renewcommand\thealgocf{\@arabic\c@algocf} % and the way it is printed
\else%                                    else
\renewcommand\thealgocf{\csname the\algocf@within\endcsname-\@arabic\c@algocf}
\fi
\renewcommand{\algocf@makecaption}[2]{%中英文双标题一定多于一行，因此去掉单行时的判断，并将\parbox中标题设置为居中
  \addtolength{\hsize}{\algomargin}%
  \sbox\@tempboxa{\algocf@captiontext{#1}{#2}}%
    \hskip .5\algomargin%
    \parbox[t]{\hsize}{\centering\algocf@captiontext{#1}{#2}}% 
  \addtolength{\hsize}{-\algomargin}%
}
\newcommand{\AlgoBiCaption}[2]{%直接取出自定义的中英文标题条目加入到真正的\caption 中  
   \caption[#1]{\protect\setlength{\baselineskip}{1.5em}#1 \protect \\ Algo. \thealgocf~~ #2} % \algoencaption{#2}   
   \addcontentsline{aen}{algoen}{\protect\numberline{\thealgoen}{#2}}
   }
%\renewcommand{\cftalgoendotsep}{1} %算法的英文目录里的点的间距

\makeatother
%定义 学科 学位
\def \xuekeEngineering {Engineering}
\def \xuekeScience {Science}
\def \xuekeManagement {Management}
\def \xuekeArts {Arts}

\ifx \xueke \xuekeEngineering
\newcommand{\cxueke}{工学}
\newcommand{\exueke}{Engineering}
\fi

\ifx \xueke \xuekeScience
\newcommand{\cxueke}{理学}
\newcommand{\exueke}{Science}
\fi

\ifx \xueke \xuekeManagement
\newcommand{\cxueke}{管理学}
\newcommand{\exueke}{Management}
\fi

\ifx \xueke \xuekeArts
\newcommand{\cxueke}{文学}
\newcommand{\exueke}{Arts}
\fi

\newcommand{\cdash}{\mbox{—\!\!\!\!—\!\!\!\!—}}%输入中文破折号的命令
\newcommand{\dif}{\mathrm{d}}%在数学模式中输入微分dx

%%% Local Variables: 
%%% mode: latex
%%% TeX-master: "../main"
%%% End: 
 % 文本格式定义
% -*-coding: utf-8 -*-

%论文版芯大小一般应为150mm×220mm（包括页眉及页码则为150mm×240mm）
%页码在版芯下边线之下隔行居中放置；
%% 左右
\setlength{\textwidth}{15cm}
\setlength{\oddsidemargin}{0.46cm}   % 左边 3cm=0.46+2.54
\setlength{\evensidemargin}{0.46cm}
%  上下
\setlength{\topmargin}{0.01cm}       % 3.3=2.54+0.76
\setlength{\headheight}{0.80cm}      % 0.8
\setlength{\headsep}{0.40cm}         % 0.4
\setlength{\textheight}{22.0cm}     % 22.0
\setlength{\footskip}{1.1cm}        %1.1

%%%%%%%%%%%%%%%%%%%%%%%%%%%%%%%%%%%%%%%%%%%%%%%%%%%%%%%%%%%
%允许公式换页显示,否则大型推导公式都在一页内，
%一页显示不下放到第二页，导致很大的空白空间，很不好看
\allowdisplaybreaks[4]

%%%%%%%%%%%%%%%%%%%%%%%%%%%%%%%%%%%%%%%%%%%%%%%%%%%%%%%%%%%
%下面这组命令使浮动对象的缺省值稍微宽松一点，从而防止幅度
%对象占据过多的文本页面，也可以防止在很大空白的浮动页上放置很小的图形。
\renewcommand{\topfraction}{0.9999999}
\renewcommand{\textfraction}{0.0000001}
\renewcommand{\floatpagefraction}{0.9999}

%%%%%%%%%%%%%%%%%%%%%%%%%%%%%%%%%%%%%%%%%%%%%%%%%%%%%%%%%%%
% 重定义一些正文相关标题
%%%%%%%%%%%%%%%%%%%%%%%%%%%%%%%%%%%%%%%%%%%%%%%%%%%%%%%%%%%
\theoremstyle{plain} \theorembodyfont{\song\rmfamily}
\theoremheaderfont{\hei\rmfamily} %\theoremseparator{:}
\newtheorem{definition}{\hei 定义}[chapter]
\newtheorem{example}{\hei 例}[chapter]
\newtheorem{algo}{\hei 算法}[chapter]
\newtheorem{theorem}{\hei 定理}[chapter]
\newtheorem{axiom}{\hei 公理}[chapter]
\newtheorem{proposition}{\hei 命题}[chapter]
\newtheorem{lemma}{\hei 引理}[chapter]
\newtheorem{corollary}{\hei 推论}[chapter]
\newtheorem{remark}{\hei 注解}[chapter]
%\newtheorem{proposition}[definition]{\hei 命题}
%\newtheorem{lemma}[definition]{\hei 引理}
%\newtheorem{exercise}[definition]{}
%\newtheorem{corollary}[definition]{\hei 推论}
%\newtheorem{remark}[definition]{\hei 注解}
%%%%%%%%%%%%%%%%%%%%%%%%%%%%%%%%%%%%%%%%%%%%%%%%%%%%%%%%%%%%%%%%%%%
%解决原proof定理环境的两个问题：
%  1. proof 中的item缩进不对
%  2. proof 中的最后一个公式下出现一个黑方块。
%\theoremsymbol{$\blacksquare$}
%\newtheorem{proof}{\hei 证明}
\newenvironment{proof}{\noindent{\hei 证明：}}{\hfill $ \square $ \vskip 4mm}
\theoremsymbol{$\square$}


%%%%%%%%%%%%%%%%%%%%%%%%%%%%%%%%%%%%%%%%%%%%%%%%%%%%%%%%%%%
% 用于中文段落缩进 和正文版式
%%%%%%%%%%%%%%%%%%%%%%%%%%%%%%%%%%%%%%%%%%%%%%%%%%%%%%%%%%%

\ifx\atempxetex\usewhat 
\newcommand{\CJKcaption}[1]{
  \ifx\CJK@actualBinding \@empty
    \PackageError{CJK}{
      You must be inside of a CJK environment to use \protect\CJKcaption}{}
  \else
    \makeatletter
    \InputIfFileExists{#1.cpx}{}{
      \PackageError{CJK}{
        Can't find #1.cpx}{
        The default captions are used if you continue.}}
    \makeatother
  \fi}  
\CJKcaption{gb_452_UTF8}
\else
\CJKcaption{gb_452_UTF8} %_UTF8
\newlength \CJKtwospaces
\def\CJKindent{
    \settowidth\CJKtwospaces{\CJKchar{"0A1}{"0A1}\CJKchar{"0A1}{"0A1}}%
    \parindent\CJKtwospaces
}
%\CJKtilde  \CJKindent
\fi

\setlength{\parindent}{26pt} %由于工大论文的每行的字距加大了，需要增加段首缩2pt

\renewcommand\contentsname{\hei 目~~~~录}

%%%%%%章节标题为“第1章”的形式
\renewcommand\chaptername{\CJKprechaptername~\thechapter~\CJKchaptername}
%%%%%%%%%%%%%%%%%%%%%%%%%%%%%%%%%%%%%%%%%%%%%%%%%%
%定义段落章节的标题和目录项的格式
%%%%%%%%%%%%%%%%%%%%%%%%%%%%%%%%%%%%%%%%%%%%%%%%%%
\setcounter{secnumdepth}{4} \setcounter{tocdepth}{2}

\titleformat{\chapter}[block]{\xiaoer\sf\hei\filcenter\boldmath}{\xiaoer\chaptertitlename}{18pt}{\xiaoer}
% \titlespacing{\chapter}{0pt}{8pt}{16pt} %block为居中对齐，hang 为悬挂的居中对齐（不包括

\titleformat{\section}[hang]{\sf\hei\xiaosan\boldmath}{\xiaosan\thesection}{0.5em}{}
% \titlespacing{\section}{0pt}{13pt}{13pt}

\titleformat{\subsection}[hang]{\sf\hei\sihao\boldmath}{\sihao\thesubsection}{0.5em}{}
% \titlespacing{\subsection}{0pt}{8pt}{7pt}

\titleformat{\subsubsection}[hang]{\sf\hei\xiaosi\boldmath}{\thesubsubsection}{0.5em}{} %[\;\;]
% \titlespacing{\subsubsection}{0pt}{3pt}{2pt}

% 章节前后的间距细调，
% 由genfazhan根据研究生院老师的审查后，与word示例范文相比细调确定
\titlespacing{\chapter}{0pt}{0mm}{6mm} %（生成的前距约12mm，后距约10mm，范例中前后距都是10mm，我调不出来，不过12mm和10mm，差不多，看不出来，应该没问题）
\titlespacing{\section}{0pt}{4mm}{4mm} %（生成的前距约7－8mm）
\titlespacing{\subsection}{0pt}{3mm}{3mm} %（生成的前距约6－7mm）
\titlespacing{\subsubsection}{0pt}{2mm}{2mm} % 

% 按工大标准, 缩小目录中各级标题之间的缩进，使它们相隔一个字符距离，也就是12pt
\makeatletter
\renewcommand*\l@chapter{\@dottedtocline{0}{0em}{4.84em}}%控制英文目录： 细点\@dottedtocline  粗点\@dottedtoclinebold
\renewcommand*\l@section{\@dottedtocline{1}{12pt}{18pt}}
\renewcommand*\l@subsection{\@dottedtocline{2}{24pt}{27pt}}
\renewcommand*\l@subsubsection{\@dottedtocline{3}{36pt}{39pt}}
\renewcommand*\l@paragraph{\@dottedtocline{4}{48pt}{48pt}}
\renewcommand*\l@subparagraph{\@dottedtocline{5}{60pt}{60pt}}
% 英文目录中章标题对应的点号及相应页码为黑体
\def\@dottedtoclinebold#1#2#3#4#5{%
 \ifnum #1>\c@tocdepth \else
 \vskip \z@ \@plus.2\p@
{\leftskip #2\relax \rightskip \@tocrmarg \parfillskip -\rightskip
\parindent #2\relax\@afterindenttrue
 \interlinepenalty\@M
\leavevmode
 \@tempdima #3\relax
 \advance\leftskip \@tempdima \null\nobreak\hskip -\leftskip
{#4}\nobreak
 \leaders\hbox{$\m@th
\mkern \@dotsep mu\hbox{\ifnum1=#1 \bf\fi.}\mkern \@dotsep
mu$}\hfill 
 \nobreak
\hb@xt@\@pnumwidth{\hfil  \normalfont \normalcolor \ifnum1=#1 \bf\fi#5}%目录层为1时，页码粗体
\par}%
\fi}
%控制中文目录
%\dottedcontents{chapter}[3.4em]{\vspace{0.5em}\hspace{-3.4em}\hei \bf\boldmath}{0.0em}{5pt}% 章标题后用粗点
\titlecontents{chapter}[3.92em]{\vspace{0.5em}\hspace{-3.92em}\hei \bf\boldmath}{\contentslabel{0em}}{\hspace*{-0em}}{\normalfont\titlerule*[5pt]{.}\contentspage} %章标题后用细点
\dottedcontents{section}[1.16cm]{}{1.8em}{5pt}
\dottedcontents{subsection}[2.00cm]{}{2.7em}{5pt}
\dottedcontents{subsubsection}[2.86cm]{}{3.4em}{5pt}

%%%%%%%%%%%%%%%%%%%%%%%%%%%%%%%%%%%%%%%%%%%%%%%%%%%%%%%
% 定义页眉和页脚 使用fancyhdr 宏包
%%%%%%%%%%%%%%%%%%%%%%%%%%%%%%%%%%%%%%%%%%%%%%%%%%%%%%%%
\newcommand{\makeheadrule}{%
\makebox[-3pt][l]{\rule[.7\baselineskip]{\headwidth}{0.4pt}}
\rule[0.85\baselineskip]{\headwidth}{2.25pt}\vskip-.8\baselineskip}
\renewcommand{\headrule}{%
    {\if@fancyplain\let\headrulewidth\plainheadrulewidth\fi
     \makeheadrule}}

\pagestyle{fancyplain}

%去掉章节标题中的数字
%%不要注销这一行，否则页眉会变成：“第1章1  绪论”样式
\renewcommand{\chaptermark}[1]{\markboth{\chaptertitlename~~ \ #1}{}}
 \fancyhf{}

%在book文件类别下,\leftmark自动存录各章之章名,\rightmark记录节标题
%% 页眉字号 工大要求 小五
%根据单双面打印设置不同的页眉；
\ifoneortwoside
  \fancyhead[RO,RE]{\song \xiaowu\leftmark}
  \fancyhead[LO,LE]{\song \xiaowu 西北工业大学\cxuewei 学位论文 }%
  \fancyfoot[C,C]{\xiaowu --~\thepage~--} %\if@mainmatter \fi
\else%
  \fancyhead[CO]{\song \xiaowu 西北工业大学\cxuewei 学位论文}
  \fancyhead[CE]{\song \xiaowu 西北工业大学\cxuewei 学位论文}%
  \fancyfoot[C,C]{\xiaowu --~\thepage~--} %\if@mainmatter   \fi
\fi

\renewcommand\frontmatter{%
    \cleardoublepage
  \@mainmatterfalse
  \pagenumbering{Roman}}
%%%%%%%%%%%%%%%%%%%%%%%%%%%%%%%%%%%%%%%%%%%%%%%%%%%%%%%%
% 设置行距和段落间垂直距离
%%%%%%%%%%%%%%%%%%%%%%%%%%%%%%%%%%%%%%%%%%%%%%%%%%%%%%%%
\ifx\CJKoglue\undefined %CJKoglue是新版CJKpunct中的命令，以此来判断采用的是新版还是旧版CJKpunct宏包
\renewcommand{\CJKglue}{\hskip 0.3pt plus 0.08\baselineskip} %加大字间距，使每行33个字 
\else
\let\CJKglue\CJKoglue
\renewcommand{\CJKglue}{\hskip 0.3pt plus 0.08\baselineskip}%加大字间距，使每行33个字
\fi
    
%\setlength{\belowcaptionskip}{10pt}   % 加大标题和表格之间的距离 \abovecaptionskip 默认是10pt
%\setlength{\parskip}{3pt plus1pt minus1pt} % 段落之间的竖直距离
\setlength{\parskip}{1.5pt plus 1pt minus 1pt} % by genfazhan（本来想改成0pt，和正常段之间一致，但生成的目录有点太紧凑了，不太好看，所以就设成了1.5pt）
\renewcommand{\baselinestretch}{1.2}% 定义行距

%%%%%%%%%%%%%%%%%%%%%%%%%%%%%%%%%%%%%%%%%%%%%%%%%%%%%%%%
% 调整列表环境的垂直间距
%%%%%%%%%%%%%%%%%%%%%%%%%%%%%%%%%%%%%%%%%%%%%%%%%%%%%%%%
\setitemize{itemindent=38pt,leftmargin=0pt,itemsep=-0.4ex,listparindent=26pt,partopsep=0pt,parsep=0.5ex,topsep=-0.25ex}
\setenumerate{itemindent=38pt,leftmargin=0pt,itemsep=-0.4ex,listparindent=26pt,partopsep=0pt,parsep=0.5ex,topsep=-0.25ex}
\setdescription{itemindent=38pt,leftmargin=0pt,itemsep=-0.4ex,listparindent=26pt,partopsep=0pt,parsep=0.5ex,topsep=-0.25ex}


%\renewcommand\@biblabel[1]{#1\hspace{0.5em}} %去除参考文献里标号两边的括号
\newcommand{\ucite}[1]{$^{\mbox{\scriptsize \cite{#1}}}$} % 增加 \ucite 命令使显示的引用为上标形式
\newcommand{\citeup}[1]{$^{\mbox{\scriptsize \cite{#1}}}$} % for WinEdt users

%%%%%%%%%%%%%%%%%%%%%%%%%%%%%%%%%%%%%%%%%%%%%%%%%%%%%%%%%%%
% 定制浮动图形和表格标题样式 %这里用ccaption完全代替了caption2的功能
\captionstyle{\centering}   %不同的图标题形式采用不同的命令
%\indentcaption{0pt}           %参见ccaption
\hangcaption
\captionnamefont{\song \rmfamily\wuhao\selectfont}
\captiontitlefont{\song \rmfamily\wuhao\selectfont}
\captiondelim{~} %~

%%%%%%%%%%%%%%%%%%%%%%%%%%%%%%%%%%%%%%%%%%%%%%%%%%%%%%%
% 定义题头格言的格式
% 用法 \begin{Aphorism}{author}
%         aphorism
%      \end{Aphorism}
\newsavebox{\AphorismAuthor}
\newenvironment{Aphorism}[1]
{\vspace{0.5cm}\begin{sloppypar} \slshape
\sbox{\AphorismAuthor}{#1}
\begin{quote}\small\itshape }
{\\ \hspace*{\fill}------\hspace{0.2cm} \usebox{\AphorismAuthor}
\end{quote}
\end{sloppypar}\vspace{0.5cm}}

%自定义一个空命令，用于注释掉文本中不需要的部分。
\newcommand{\comment}[1]{}

\renewcommand\contentsname{\hei 目~~~~录}
\renewcommand\listfigurename{\hei 插~~~~图}
\renewcommand\listtablename{\hei 表~~~~格}

%%%%%%将章标题中的中文数字（一、二、三）变为阿拉伯数字(1,2,3)
\renewcommand\CJKthechapter{%\CJKnumber
{\@arabic\c@chapter}}

%%%%%%不要拉大行距使得页面充满
\raggedbottom

% This is the flag for longer version
\newcommand{\longer}[2]{#1}

\newcommand{\ds}{\displaystyle}

% define graph scale
\def\gs{1.0}

%%%%%%%%%%%%%%%%%%%%%%%%%%%%%%%%%%%%%%%%%%%%%%%%%%%%%%%%%%%%%%%%%%%%%%
% 自定义项目列表标签及格式 \begin{hitlist} 列表项 \end{hitlist}
%%%%%%%%%%%%%%%%%%%%%%%%%%%%%%%%%%%%%%%%%%%%%%%%%%%%%%%%%%%%%%%%%%%%%%
\newcounter{hitctr} %自定义新计数器
\newenvironment{hitlist}{%%%%%定义新环境
\begin{list}{{\hei (\arabic{hitctr})}} %%标签格式
    {
     \usecounter{hitctr}
     \setlength{\leftmargin}{0cm}     %左边界
     \setlength{\parsep}{0ex}         %段落间距
     \setlength{\topsep}{0pt}         %列表到上下文的垂直距离
     \setlength{\itemsep}{0ex}        %标签间距
     \setlength{\labelsep}{0.3em}     %标号和列表项之间的距离,默认0.5em
     \setlength{\itemindent}{46pt}    %标签缩进量
     \setlength{\listparindent}{27pt} %段落缩进量
    }}
{\end{list}}%%%%%

%%%%%%%%%%%%%%%%%%%%%%%%%%%%%%%%%%%%%%%%%%%%%%%%%%%%%%%%%%%%%%%%%%%%%%
% 自定义项目列表标签及格式 \begin{publist} 列表项 \end{publist}
%%%%%%%%%%%%%%%%%%%%%%%%%%%%%%%%%%%%%%%%%%%%%%%%%%%%%%%%%%%%%%%%%%%%%%
\newcounter{pubctr} %自定义新计数器
\newenvironment{publist}{%%%%%定义新环境
\begin{list}{\arabic{pubctr}} %%标签格式
    {
     \usecounter{pubctr}
     \setlength{\leftmargin}{2em}     % 左边界 \leftmargin =\itemindent + \labelwidth + \labelsep
     \setlength{\itemindent}{0em}     % 标号缩进量
     \setlength{\labelwidth}{1em}     % 标号宽度
     \setlength{\labelsep}{1em}       % 标号和列表项之间的距离,默认0.5em
     \setlength{\rightmargin}{0em}    % 右边界
     \setlength{\topsep}{0ex}         % 列表到上下文的垂直距离
%     \setlength{\partopsep}{0ex}      % 列表是一个新的段落时，加的额外到上下文的距离
     \setlength{\parsep}{0ex}         % 段落间距
     \setlength{\itemsep}{0ex}        % 标签间距
     \setlength{\listparindent}{26pt} % 段落缩进量
    }}
{\end{list}}%%%%%

%%%%%%%%%%%%%%%%%%%%%%%%%%%%%%%%%%%%%%%%%%%%%%%%%%%%%%%%%%%%%%%%%%%%%%
% 默认字体
\renewcommand\normalsize{%
  \@setfontsize\normalsize{12pt}{13pt}
  %\setlength\abovedisplayskip{8pt plus 2pt minus 2pt}
  %\setlength\abovedisplayshortskip{7pt plus 2pt minus 2pt}  
  \setlength\abovedisplayskip{5pt plus 2pt minus 2pt} %by gengfazhan 
  \setlength\abovedisplayshortskip{4pt plus 2pt minus 2pt} %by gengfazhan（生成的公式和正文距离减小了，比正常行距稍大一些）
  \setlength\belowdisplayskip{\abovedisplayskip}
  \setlength\belowdisplayshortskip{\abovedisplayshortskip}
  \setlength\jot{6pt}
  \let\@listi\@listI}
\def\defaultfont{\renewcommand{\baselinestretch}{1.5}\normalsize\selectfont} %正文行间距，大致与word中小四宋体的1.25倍行距相同，使用者若不满意可精调此参数
\predisplaypenalty=0  %公式之前可以换页，公式出现在页面顶部

%%%%%%%%%%%%%%%%%%%%%%%%%%%%%%%%%%%%%%%%%%%%%%%%%%%%%%%%%%%%%%%%%%%%%%
% 封面、摘要、版权、致谢格式定义
\def\ctitle#1{\def\@ctitle{#1}}\def\@ctitle{}
\def\cdegree#1{\def\@cdegree{#1}}\def\@cdegree{}
\def\caffil#1{\def\@caffil{#1}}\def\@caffil{}
\def\csubject#1{\def\@csubject{#1}}\def\@csubject{}
\def\cauthor#1{\def\@cauthor{#1}}\def\@cauthor{}
\def\csupervisor#1{\def\@csupervisor{#1}}\def\@csupervisor{}
\def\cassosupervisor#1{\def\@cassosupervisor{~ & {\hei 副 \hfill 导 \hfill 师：} & #1\\}}\def\@cassosupervisor{}
\def\ccosupervisor#1{\def\@ccosupervisor{~ & {\hei 联 \hfill 合\hfill 导 \hfill 师：} & #1\\}}\def\@ccosupervisor{}
\def\cdate#1{\def\@cdate{#1}}\def\@cdate{}
\long\def\cabstract#1{\long\def\@cabstract{#1}}\long\def\@cabstract{}
\def\ckeywords#1{\def\@ckeywords{#1}}\def\@ckeywords{}

\def\etitle#1{\def\@etitle{#1}}\def\@etitle{}
\def\edegree#1{\def\@edegree{#1}}\def\@edegree{}
\def\eaffil#1{\def\@eaffil{#1}}\def\@eaffil{}
\def\esubject#1{\def\@esubject{#1}}\def\@esubject{}
\def\eauthor#1{\def\@eauthor{#1}}\def\@eauthor{}
\def\esupervisor#1{\def\@esupervisor{#1}}\def\@esupervisor{}
%\def\eassosupervisor#1{\def\@eassosupervisor{#1}}\def\@eassosupervisor{}
\def\eassosupervisor#1{\def\@eassosupervisor{~ & \textbf{Associate Supervisor:} & #1\\}}\def\@eassosupervisor{}
%\def\ecosupervisor#1{\def\@ecosupervisor{#1}}\def\@ecosupervisor{}
\def\ecosupervisor#1{\def\@ecosupervisor{~ & \textbf{Co Supervisor:} & #1\\}}\def\@ecosupervisor{}
\def\edate#1{\def\@edate{#1}}\def\@edate{}
\long\def\eabstract#1{\long\def\@eabstract{#1}}\long\def\@eabstract{}
\long\def\NotationList#1{\long\def\@NotationList{#1}}\long\def\@NotationList{}
\def\ekeywords#1{\def\@ekeywords{#1}}\def\@ekeywords{}
\def\natclassifiedindex#1{\def\@natclassifiedindex{#1}}\def\@natclassifiedindex{}
\def\internatclassifiedindex#1{\def\@internatclassifiedindex{#1}}\def\@internatclassifiedindex{}
\def\statesecrets#1{\def\@statesecrets{#1}}\def\@statesecrets{}

%%%%%%%%%%%%%%%%%%%%%%%%%%%%%%%%%%%%%%%%%%%%%%%%%%%%%%%%%%%%%%%
% 定义封面
\def\makecover{
    \normalbiao %表格字号设置
    \begin{titlepage}
    
    % 封面一
    %\begin{center}

    %\parbox[t][1cm][b]{\textwidth}{\xiaoyi
    %\begin{center} {\song \cxuewei 学位论文 }\end{center} } %\ifxueweimaster\cxueke\fi

    %\parbox[t][0.8cm][t]{\textwidth}{
    %\begin{center} \end{center} }

    %\parbox[t][2cm][t]{\textwidth}{\erhao
    %\begin{center} {\hei  \@ctitle}\end{center} }

    %\ifxueweidoctor
    %\parbox[t][3.8cm][t]{\textwidth}{\erhao %英文标题太长时可以采用\xiaoer
    %\begin{center} {\bfseries  \@etitle}\end{center} }
    %\else
    %\parbox[t][3.8cm][t]{\textwidth}{\centering
        %\ }
    %\fi

    %\parbox[t][2.0cm][t]{\textwidth}{\xiaoer
    %\begin{center} {\song  \@cauthor}  \end{center} }

    %\ifxueweidoctor
    %\parbox[t][8.5cm][t]{\textwidth}{\centering
        %\includegraphics[width = 8cm]{hit_logo}}
    %\else
    %\parbox[t][8.5cm][t]{\textwidth}{\centering
        %\includegraphics[width = 8cm]{hit_logo} }
    %\fi

    %\parbox[t][1.2cm][t]{\textwidth}{\xiaoer
    %\begin{center} {\kai  西北工业大学}  \end{center} }

    %\parbox[t][0.5cm][t]{\textwidth}{
    %\begin{center} {\song \xiaoer \@cdate} \end{center} }

    %\end{center}

    % 封二 空白页
    %\ifoneortwoside
      %\newpage
      %~~~\vspace{1em}
      %\thispagestyle{empty}
    %\fi

    %内封
    \newpage
    \thispagestyle{empty}
    \begin{center}
    \parbox[t][0.6cm][t]{\textwidth}{
    \begin{center} \end{center}}

    %\parbox[t][2.2cm][t]{\textwidth}{
    %\song \xiaosi 国内图书分类号： \@natclassifiedindex \\
                 %国际图书分类号： \@internatclassifiedindex }
    %\parbox[t][2.2cm][t]{\textwidth}{\song \xiaosi \centering 
    %\begin{tabular}{>{\raggedleft}p{4cm}>{\raggedright}p{3.5cm}>{\raggedleft}p{3.5cm}>{\raggedright}p{2cm}}%% %% {rlrl}
        %分类号：& \@natclassifiedindex  & 学校代码：& 10699 \tabularnewline 
        %密~~~~级：&  公开 & 学~~~~~~~~号：& 067100540
    %\end{tabular}}

    %\parbox[t][2.7cm][b]{\textwidth}{\xiaoer
    %\begin{center} {\song  \@cdegree 学位论文 }\end{center} }

    \setlength{\baselineskip}{1.5\baselineskip}
    \parbox[t][3.0cm][b]{\textwidth}{\yihao
    \begin{center} {\hei  \@ctitle}\end{center} }

    \parbox[t][4.3cm][t]{\textwidth}{
    \begin{center}\xiaoer (申请西北工业大学工学博士学位论文) \end{center} }

    \parbox[t][10cm][c]{\textwidth}{ {\sanhao
    \begin{center} \renewcommand{\arraystretch}{1.5} \song 
    \begin{tabular}{lll@{\extracolsep{0em}}l}
    ~ & {\hei 作 \hfill  者：}                            & \@cauthor\\
    ~ & {\hei 指\hfill 导\hspace{1em}教\hfill 师：}       & \@csupervisor\\
    %\@ccosupervisor
    %\@cassosupervisor
    %~ & {\hei 申\hfill 请\hspace{1em}学\hfill 位：} & \@cdegree\\
    ~ & {\hei 学\hfill 科\hspace{1em}专\hfill 业：}       & \@csubject\\
    ~ & {\hei 培\hfill 养\hspace{1em}单\hfill 位：}       & \@caffil\\
    ~ & {\hei 申\hfill 请\hspace{1em}日\hfill 期：}       & \@cdate\\
    \end{tabular} \renewcommand{\arraystretch}{1}
    \end{center} } }
\end{center}

%%%%%%增加一空白页
  \ifoneortwoside
    \newpage
    ~~~\vspace{1em}
    \thispagestyle{empty}
  \fi

    % 英文封面
    \newpage
    \thispagestyle{empty}
    \begin{center}
    \parbox[t][0.6cm][t]{\textwidth}{
    \begin{center} \end{center}}

    %\parbox[t][2.2cm][t]{\textwidth}{
    %\xiaosi Classified Index: \@natclassifiedindex \\
                  %U.D.C.:  \@internatclassifiedindex
                  %}

    %\parbox[t][2.7cm][b]{\textwidth}{\xiaoer
    %\begin{center} {  Dissertation for the {\exueweier} Degree }\end{center} } %与中文保持一致，删除in {\exueke}

    \parbox[t][1.0cm][b]{\textwidth}{\xiaoyi
    \begin{center} {\bf \@etitle}\end{center} }

    \parbox[t][2.0cm]{\textwidth}{ 
    \begin{center} 
        {\sanhao
        \vspace{3cm}
        By \@eauthor\\
        \vspace{.3cm}
        Supervisor:  \@esupervisor\\

        \vspace{3cm}
        A Dissertation Submitted to the \\
        \vspace{.3cm}
        {\bf GRADUATE SCHOOL OF NORTHWESTERN \\ 
        \vspace{.3cm}
        POLYTECHNICAL UNIVERSITY \\}

        \vspace{2.2cm}
        In Partial Fulfillment of the Requirements for the Degree \\
        \vspace{.3cm}
        {\bf DOCTORAL OF PHILOSOPHY} \\
        \vspace{.3cm}
        (\@esubject) \\

        \vspace{3cm}
        \@edate \\
        }
    \end{center} }

    %\parbox[t][6cm][c]{\textwidth}{ {\sihao
    %\begin{center} \renewcommand{\arraystretch}{1.5} 
    %\begin{tabular}{p{0cm}p{14em}p{13.4em}l@{\extracolsep{2em}}l}
    %~ & \textbf{Candidate:}                     &  \@eauthor\\
    %~ & \textbf{Supervisor:}                    &  \@esupervisor\\
    %\@ecosupervisor
    %\@eassosupervisor
    %~ & \textbf{Academic Degree Applied for:}   &  \@edegree\\
    %~ & \textbf{Specialty:}                     &  \@esubject\\
    %~ & \textbf{Affiliation:}                   &  \@eaffil\\
    %~ & \textbf{Date:}               &  \@edate\\
    %~ & \textbf{Degree-Conferring-Institution:} &  Northwestern Polytechnical University
    %\end{tabular}
    %\end{center}}}

    \end{center}
    \end{titlepage}

%%%%%%增加一空白页
  \ifoneortwoside
    \newpage
    ~~~\vspace{1em}
    \thispagestyle{empty}
  \fi
%%%%%%%%%%%%%%%%%%%   Abstract and keywords  %%%%%%%%%%%%%%%%%%%%%%%
\clearpage \BiAppendixChapter{摘~~~~要}{Abstract (in Chinese)} %不要挪到下一行，生成正确的摘要toc
\setcounter{page}{1}
\song \normalsize
\defaultfont
\@cabstract
\vspace{1em}

\hangafter1\hangindent4.28em\noindent
{\hei 关键词：} \@ckeywords %新规范有了冒号之后，不用再加\quad 空白距离了。

%%%%%%%%%%%%%%%%%%%   English Abstract  %%%%%%%%%%%%%%%%%%%%%%%%%%%%%%
\clearpage
\defaultfont \BiAppendixChapter{\textbf{Abstract}}{Abstract (in English)} %不要挪到下一行，生成正确的摘要toc
\@eabstract
\vspace{1em}

\hangafter1\hangindent5.5em\noindent
{\textbf{Keywords:}}  \@ekeywords % 新规范有了冒号之后，不用再加\quad 空白距离了。
\wuhaobiao  %正文表格设置
}  %\makecover

%%%%%%%%%%%%%%%%%%%%%%%%%%%%%%%%%%%%%%%%%%%%%%%%%%%%%%%%%%%%%%%
% 英文目录格式
\def\@dotsep{1}           % 定义英文目录的点间距
\setlength\leftmargini {0pt}
\setlength\leftmarginii {0pt}
\setlength\leftmarginiii {0pt}
\setlength\leftmarginiv {0pt}
\setlength\leftmarginv {0pt}
\setlength\leftmarginvi {0pt}

\def\engcontentsname{\bfseries Contents}
\newcommand\tableofengcontents{%
   %\cleardoublepage
   \pdfbookmark[0]{Contents}{econtent}
   \if@twocolumn
     \@restonecoltrue\onecolumn
   \else
     \@restonecolfalse
   \fi   \chapter*{\engcontentsname  %chapter*上移一行，避免在toc中出现。
       \@mkboth{%
          \engcontentsname}{\engcontentsname}}
   \@starttoc{toe}%
   \if@restonecol\twocolumn\fi
   }

\urlstyle{same}  %论文中引用的网址的字体默认与正文中字体不一致，这里修正为一致的。

%主要符号表 \NotationList
\long\def\notation{\clearpage \BiAppendixChapter{主要符号表}{Main Symbol Table}\normalbiao\@NotationList\wuhaobiao}

%%% 五号字表格设置 start %%%%%
\gdef\tpltable{\relax}
\let\tpltable\longtable
\gdef\wuhaobiao{%五号字
    \def\tabular{\wuhao\gdef\@halignto{}\@tabular}
    \def\endtabular{\endarray $\egroup}
    \def\longtable{\wuhao\tpltable}
    \def\endlongtable{\adl@LTlastrow \adl@org@endlongtable}
}
\gdef\normalbiao{%正常字号
    \def\tabular{\gdef\@halignto{}\@tabular}
    \def\endtabular{\endarray $\egroup}
    \def\longtable{\tpltable}
    \def\endlongtable{\adl@LTlastrow \adl@org@endlongtable}
}
\wuhaobiao
%%% 五号字表格设置 end %%%%%
%\renewcommand{\arraystretch}{1.4} %表格中行距 ，导致公式 bmatrix 间距增大。

% 表格与下方间距
\renewcommand\endtable{\vspace{-4mm}\end@float}
% 算法与下方间距
%\renewcommand\endalgorithm{\@algocf@finish \ifthenelse {\equal {\algocf@float }{figure}}{\end {figure}}{
%\@algocf@term@caption \ifthenelse {\boolean {algocf@algoH}}{\end {algocf@Here}}
%{\end {algocf}}}\@algocf@term\vspace{-5mm}}
% algorithm2e的4.0版本删除了\endalgorithm的命令，改为直接定义algoco@algorithm环境的形式，造成使用上面的语句为出现错误：
% Latex Error:\begin{algoco@algorithm} on input line 526 ended by \end{algorithm}  
\makeatother


%%%%%%%%%%%%%%%%%%%%%%%%%%%%%%%%%%%%%%%%%%%%%%%%%%%%%%%%%%%
% 正文部分
%%%%%%%%%%%%%%%%%%%%%%%%%%%%%%%%%%%%%%%%%%%%%%%%%%%%%%%%%%%
\frontmatter
\sloppy % 解决中英文混排的断行问题，会加入间距，但不会影响断行
% -*-coding: utf-8 -*-

\newcommand{\chinesethesistitle}{西北工业大学硕博士学位论文~\LaTeX~模板} %授权书用，无需断行
\newcommand{\englishthesistitle}{\LaTeX~Dissertation Template of Northwestern Polytechnical University} %\uppercase作用：将英文标题字母全部大写；
\newcommand{\chinesethesistime}{2011~年~11~月}  %封面底部的日期中文形式
\newcommand{\englishthesistime}{November 2011}    %封面底部的日期英文形式

\ctitle{西北工业大学硕博士学位论文\\~\LaTeX~模板}  %封面用论文标题，自己可手动断行
\cdegree{\cxueke\cxuewei}
\csubject{计算机科学与技术}                 %(~按二级学科填写~)
\caffil{西北工业大学} %（在校生填所在系名称，同等学力人员填工作单位）
\cauthor{某~~某~~某}
\csupervisor{某~~某~~某~~~~教~~授} %导师名字
%\cassosupervisor{某~~~~~~某~~~~教~~授}     %(~如无副导师可以不列此项~)
%\ccosupervisor{某~~某~~某~~~~教~~授~} %(~如无联合培养导师则不列此项~)
\cdate{\chinesethesistime}

\etitle{\englishthesistitle}
\edegree{\exuewei \ of \exueke}
\esubject{Computer Science and  Technology}  %英文二级学科名
\eaffil{School \hfill of\hfill Computer\hfill Science}%英文单位 %换行用\newline，不要用\\
\eauthor{Alice}                   %作者姓名 （英文）
\esupervisor{Prof. Lambda}       % 导师姓名 （英文）
%\ecosupervisor{Professor X}
%\eassosupervisor{Professor Y}
\edate{\englishthesistime}

\natclassifiedindex{TP393}  %国内图书分类号
\internatclassifiedindex{681.324}  %国际图书分类号
\statesecrets{公开} %秘密

\iffalse
\BiAppendixChapter{摘~~~~要}{}  %使用winedt编辑时文档结构图（toc）中为了显示摘要，故增加此句；
\fi
\cabstract{
该模板是以哈尔滨工业大学\LaTeX硕博士学位论文模板为基础(哈工大
PlutoThesis项目，UTF8版本，svn版本号为r132)，根据西北工业大学学位论文
规范制作的\LaTeX硕博士论文模板。该模板使用pdflatex在Linux环境下
(Ubuntu904, textlive 2009)通过编译测试。

模板主要做了以下几方面的修改：
\begin{itemize}[,leftmargin=4em,itemindent=0em]
\renewcommand{\labelitemi}{{\large$\bullet$}}
\item 移除了英文图表和目录，修改了封面的内容和页眉的显示方式，添加了知
    识产权声明，修改了参考文献的格式。
    
\item 做了些小的调整，具体修改请diff r132。

\end{itemize} 
}

\ckeywords{\LaTeX，论文模板，西北工业大学}


\eabstract{
This is a \LaTeX dissertation template of Northwestern Polytechnical
University, which is built according to the required format. The
template is derived from PlutoThesis project.
}

\ekeywords{\LaTeX, dissertation template, Northwestern Polytechnical
University}

\makecover
\clearpage
 % 封面

%% 中英目录
\renewcommand{\baselinestretch}{1}
\fontsize{12pt}{12pt}\selectfont
\clearpage{\pagestyle{empty}\cleardoublepage}
\pdfbookmark[0]{目~~~~录}{mulu}
\tableofcontents    % 中文目录
\ifxueweidoctor     % 英文目录右开
  \clearpage{\pagestyle{empty}\cleardoublepage}
\else%
  \ifoneortwoside\clearpage{\pagestyle{empty}\cleardoublepage}\fi
\fi

%% -*-coding: utf-8 -*-

%% 中英文插图、表格、算法索引   
%% 硕博士学位论文规范均不要求这一项，请正式打印的时候在main.tex中屏蔽掉% -*-coding: utf-8 -*-

%% 中英文插图、表格、算法索引   
%% 硕博士学位论文规范均不要求这一项，请正式打印的时候在main.tex中屏蔽掉% -*-coding: utf-8 -*-

%% 中英文插图、表格、算法索引   
%% 硕博士学位论文规范均不要求这一项，请正式打印的时候在main.tex中屏蔽掉\input{figtab.tex}；

\ifxueweidoctor
  %\clearpage{\pagestyle{empty}\cleardoublepage}   % 清除目录后面空页的页眉和页脚
\else%
  {\ifoneortwoside\clearpage{\pagestyle{empty}\cleardoublepage}\else\newpage\fi} % 清除目录后面空页的页眉和页脚
\fi

\addcontentsline{toc}{chapter}{\hei 插~~~~图}   % 中文插图加入到中文目录
\listoffigures                                  % 生成中文 图形索引

\ifxueweidoctor
  %\clearpage{\pagestyle{empty}\cleardoublepage}   % 清除目录后面空页的页眉和页脚
\else%
  {\ifoneortwoside\clearpage{\pagestyle{empty}\cleardoublepage}\else\newpage\fi}   % 清除目录后面空页的页眉和页脚
\fi

\addcontentsline{toc}{chapter}{\hei 表~~~~格}   % 中文表格加入到中文目录
\listoftables                                   % 生成中文 表格索引


\ifxueweidoctor
  %\clearpage{\pagestyle{empty}\cleardoublepage}   % 清除目录后面空页的页眉和页脚
\else%
  {\ifoneortwoside\clearpage{\pagestyle{empty}\cleardoublepage}\else\newpage\fi}   % 清除目录后面空页的页眉和页脚
\fi

%%% Local Variables: 
%%% mode: latex
%%% TeX-master: "../main"
%%% End: 
；

\ifxueweidoctor
  %\clearpage{\pagestyle{empty}\cleardoublepage}   % 清除目录后面空页的页眉和页脚
\else%
  {\ifoneortwoside\clearpage{\pagestyle{empty}\cleardoublepage}\else\newpage\fi} % 清除目录后面空页的页眉和页脚
\fi

\addcontentsline{toc}{chapter}{\hei 插~~~~图}   % 中文插图加入到中文目录
\listoffigures                                  % 生成中文 图形索引

\ifxueweidoctor
  %\clearpage{\pagestyle{empty}\cleardoublepage}   % 清除目录后面空页的页眉和页脚
\else%
  {\ifoneortwoside\clearpage{\pagestyle{empty}\cleardoublepage}\else\newpage\fi}   % 清除目录后面空页的页眉和页脚
\fi

\addcontentsline{toc}{chapter}{\hei 表~~~~格}   % 中文表格加入到中文目录
\listoftables                                   % 生成中文 表格索引


\ifxueweidoctor
  %\clearpage{\pagestyle{empty}\cleardoublepage}   % 清除目录后面空页的页眉和页脚
\else%
  {\ifoneortwoside\clearpage{\pagestyle{empty}\cleardoublepage}\else\newpage\fi}   % 清除目录后面空页的页眉和页脚
\fi

%%% Local Variables: 
%%% mode: latex
%%% TeX-master: "../main"
%%% End: 
；

\ifxueweidoctor
  %\clearpage{\pagestyle{empty}\cleardoublepage}   % 清除目录后面空页的页眉和页脚
\else%
  {\ifoneortwoside\clearpage{\pagestyle{empty}\cleardoublepage}\else\newpage\fi} % 清除目录后面空页的页眉和页脚
\fi

\addcontentsline{toc}{chapter}{\hei 插~~~~图}   % 中文插图加入到中文目录
\listoffigures                                  % 生成中文 图形索引

\ifxueweidoctor
  %\clearpage{\pagestyle{empty}\cleardoublepage}   % 清除目录后面空页的页眉和页脚
\else%
  {\ifoneortwoside\clearpage{\pagestyle{empty}\cleardoublepage}\else\newpage\fi}   % 清除目录后面空页的页眉和页脚
\fi

\addcontentsline{toc}{chapter}{\hei 表~~~~格}   % 中文表格加入到中文目录
\listoftables                                   % 生成中文 表格索引


\ifxueweidoctor
  %\clearpage{\pagestyle{empty}\cleardoublepage}   % 清除目录后面空页的页眉和页脚
\else%
  {\ifoneortwoside\clearpage{\pagestyle{empty}\cleardoublepage}\else\newpage\fi}   % 清除目录后面空页的页眉和页脚
\fi

%%% Local Variables: 
%%% mode: latex
%%% TeX-master: "../main"
%%% End: 
  %图表索引, 如果不需要图表索引，注释掉这一句即可；
% \notation  %主要符号表
\addtocontents{toc}{\protect\vskip1\baselineskip} % 中文目录增加空行
\addtocontents{toe}{\protect\vskip1\baselineskip} % 英文目录增加空行

\ifxueweidoctor
  \clearpage{\pagestyle{empty}\cleardoublepage}   % 清除目录后面空页的页眉和页脚
\else%
  \ifoneortwoside\clearpage{\pagestyle{empty}\cleardoublepage}\fi  % 清除目录后面空页的页眉和页脚
\fi                                               %  第一章是否右开

\mainmatter
\defaultfont % 对应于小四的标准字号为12pt, 可以在正文中用此命令修改所需要字体的的大小

% -*-coding: utf-8 -*-

\defaultfont

\BiChapter{ 绪~~~论}{ }

\BiSection{背景 }{ }

虽然微软Office所见即所得的方式很容易上手，但当用于较为正式的写作时还是
会让作者花费大量时间在论文本身之外的排版工作。\LaTeX可靠而稳定的排版系
统可以让作者将精力集中于论文写作本身。如果你开始尝试使用\LaTeX进行论文
或者其他较为正式的文稿写作，你会慢慢爱上她。

参考文献对应的BibTeX文件位于\texttt{reference/reference.bib}，这里推荐
使用JabRef\footnote{\url{http://jabref.sourceforge.net/}}对参考文献进
行管理。中英文摘要位于\texttt{preface}目录下的\texttt{cover.tex}文件中
。\texttt{appendix}目录中包含了致谢、论文发表情况、科研经历和知识产权
声明。\texttt{body}目录中包含了论文的主要章节。\texttt{setup}目录则是
一些相关的设置。


表~\ref{table:intro-motes} 给出了一个表的例子。图片一般放在
\texttt{figures}目录下，图~\ref{fig:intro-npu} 给出了一个图片的例子。
文献~\cite{lukaicheng2002} 是一本书，文献~\cite{wsn:jos} 是一篇期刊的
论文。

\begin{table}[htbp]
\TableBiCaption{ 几种常见传感器节点硬件基本参数 } {}
\vspace{2mm}
\label{table:intro-motes}
\centering
\begin{tabular}{ l l l l l l }
\toprule
\bfseries            &\bfseries              &\bfseries Data       &\bfseries Program    &\bfseries External   &\bfseries Data\\
\bfseries            &\bfseries Processor    &\bfseries memory     &\bfseries memory     &\bfseries flash      &\bfseries rate   \\
\bfseries Platforms  &\bfseries (MHz / bits) &\bfseries (KB)       &\bfseries (KB)       &\bfseries (KB)       &\bfseries (kbps) \\
\midrule
              Mica2  & 7.4 / 8      & 4          & 128    & 512    & 38 \\
              MicaZ  & 7.4 / 8      & 4          & 128    & 512    & 250 \\
              TelosB & 4.0 / 16     & 10         &  48    & 1024   & 250 \\
              Iris   & 12.0 / 8     & 8          & 128    & 512    & 250 \\
\bottomrule
\end{tabular}
\end{table}

\begin{figure}[htbp]
\centering
\includegraphics[width = 0.9\textwidth]{npu}
\FigureBiCaption{西北工业大学长安校区}{}
\label{fig:intro-npu}
\end{figure}

\BiSection{其他说明}{ }

论文使用pdflatex编译，在Linux系统下如果需要使用Windows下的字体，可能需
需要制作字体，具体方法请上网Google。另外，也可尝试对字体支持更好的
xetex编译，但可能需要进行一些修改。

main.tex中有每一章从右侧/奇数页开始的选项(openleft)，但插入的并不是一
个完全的空白页。如需空白页，可以关闭openleft选项，然后手动在合适位置使
用如下命令插入空白页：
\begin{verbatim}
\newpage
~~~\vspace{1em}
\thispagestyle{empty}
\end{verbatim}

编译时直接使用\texttt{make}命令，得到\texttt{main\_pdflatex.pdf}文件。


% -*-coding: utf-8 -*-

\defaultfont

\BiChapter{ 总结与展望 }{ }

\newpage
~~~\vspace{1em}
\thispagestyle{empty}




%参考文献
\defaultfont
\ifx\atempxetex\usewhat
\bibliographystyle{chinesebst2005_UTF8}
\else
\bibliographystyle{chinesebst2005_UTF8}
\fi
\addcontentsline{toc}{chapter}{\hei \ReferenceCName}      % 参考文献加入到中文目录
\addcontentsline{toe}{chapter}{\bfseries \xiaosi \ReferenceEName} % 参考文献加入到英文目录
\addtolength{\bibsep}{-0.8 em} \nocite{*}
\bibliography{reference/reference}

\ifoneortwoside
\newpage
~~~\vspace{1em}
\thispagestyle{empty}
\fi

%\addtocontents{fen}{\protect\vskip1\baselineskip}
%\addtocontents{ten}{\protect\vskip1\baselineskip}
%英文图形和表格索引里加入空白行，通常放在 \include{appendix/appA}% 附录A之前。
%区分开正文和附录的图形和表格，一般没有这个必要。

% -*-coding: utf-8 -*-

\defaultfont
\BiAppendixChapter{攻读\cxuewei 学位期间发表的学术论文} {Papers
Published in the Period of Ph. D. Education}

% 参见...\Accessories\ThesisCriterion\参考文献国标GB7714-2005.pdf内规定
（一）国际会议论文
\begin{publist}
\item \underline{San Zhang}, Si Li, Er Wang. title, \textit{conference
    name}, publisher, city, Country, year. (\textbf{EI Compendex,
    Access Number: xxxxxxxxxxxxxx})

\end{publist}

（二）期刊论文
\begin{publist}
\item \underline{张三}, 李四. 题目, 期刊, 年, 卷(期): 页码
\end{publist}
 

%（二）申请及已获得的专利（无专利时此项可以不列出）


%（三）获得的科技奖励（无获奖时此项不必列出）

% 注：
% 如已发表的学术论文被EI或SCI收录，请标明收录号及SCI论文的影响因子；
% 对已接收但尚未发表出来的学术论文，请注明是否EI或SCI刊源。

\ifoneortwoside
\newpage
~~~\vspace{1em}
\thispagestyle{empty}
\fi



    % 所发文章
% -*-coding: utf-8 -*-

\defaultfont

\BiAppendixChapter{攻读\cxuewei 学位期间科研工作和获奖情况}{Resume}

{\hei 科研工作}
\begin{publist}
\item “项目名称”，国家自然科学基金课题 (项目编号)
\item “项目名称”，国家科技支撑计划 (项目编号)
\end{publist}

%{\hei 获奖情况}
%\begin{publist}
%\item xxxx~年~x~月~~~xxxxxxx
%\end{publist}

\vspace{1cm}
{\hei 专利情况}
\begin{publist}
\item 张三, 李四, \underline{王二}. 专利名称. 发明专利 (专利号)
\end{publist}

\ifoneortwoside
\newpage
~~~\vspace{1em}
\thispagestyle{empty}
\fi


          % 个人简历
% -*-coding: utf-8 -*-

\defaultfont

\BiAppendixChapter{致~~~~谢}{Acknowledgement}

\ifoneortwoside
\newpage
~~~\vspace{1em}
\thispagestyle{empty}
\fi



% 致谢
% -*-coding: utf-8 -*-

\newpage
\thispagestyle{empty}

\begin{center}
    {\xiaoer {\hei 西北工业大学 \\ 学位论文知识产权声明书}}
\end{center}

本人完全了解学校有关保护知识产权的规定，即：研究生在校攻读学位期间论文
工作的知识产权单位属于西北工业大学。学校有权保留并向国家有关部门或机构
送交论文的复印件和电子版。本人允许论文被查阅和借阅。学校可以将本学位论
文的全部或部分内容编入有关数据库进行检索，可以采用影印、缩印或扫描等复
制手段保存和汇编本学位论文。同时本人保证，毕业后结合学位论文研究课题再
撰写的文章一律注明作者单位为西北工业大学。

保密论文待解密后适用本声明。

\begin{flushleft}
~~~~~~~~学位论文作者签名：~~~~~~~~~~~~~~~~~~~~~~~~~~~~~指导教师签名:~~~~~~~~~~~~~~~~~~~~~~~~~~~~~~~~~~~~~~~~~

~~~~~~~~~~~~~~~~~~~~~~~~~~~~~~~~~~~~~~~~~年~~~~~~月~~~~~~日~~~~~~~~~~~~~~~~~~~~~~~~~~~~~~~~~年~~~~~~月~~~~~~日~~~~
\end{flushleft}

\vspace{1cm}
\begin{flushleft}
----------------------------------------------------------------------------------------------------------
\end{flushleft}
\vspace{1cm}

\begin{center}
    {\xiaoer {\hei 西北工业大学 \\ 学位论文原创性声明}}
\end{center}

秉承学校严谨的学风和优良的科学道德,本人郑重声明:所呈交的学位论文,是本
人在导师的指导下进行研究工作所取得的成果。尽我所知,除文中已经注明引用
的内容和致谢的地方外,本论文不包含任何其他个人或集体已经公开发表或撰写
过的研究成果,不包含本人或其他已申请学位或其他用途使用过的成果。对本文
的研究做出重要贡献的个人和集体,均已在文中以明确方式表明。

本人学位论文与资料若有不实,愿意承担一切相关的法律责任。

\begin{flushright}

~~~~~~~~~~~~~~~~~~~~~~~~~~~~~学位论文作者签名：~~~~~~~~~~~~~~~~~~~~~~~~~~~~~~\\          
~~~~~~~~~~~~~~~~~~~~~~~~~~~~~~~~~~~~~~~~~年~~~~~~月~~~~~~日~~~~~~~~~~
\end{flushright}

   % 承诺

\clearpage
\ifx\atempxetex\usewhat\else
\end{CJK*}
\fi

\end{document}

%%% Local Variables: 
%%% mode: latex
%%% TeX-master: t
%%% End: 
